\section{Products of cubic fourfolds with geometric K3 categories}
\label{sec:cubic-products}

The assertions in this section are conditional on the specified theorem in the
OpenAI release.  The argument supplies a corollary of that theorem through an
established motivic comparison; it does not independently verify the released
proof or establish priority for the corollary.

\begin{assumption}[K3-product input]
\label{ass:k3-products}
Theorem~1.1 of \cite{k3} holds: for every finite list of smooth projective
complex K3 surfaces $S_1,\ldots,S_m$, allowing repeated factors, and every
integer $p$, the rational cycle-class map
\[
 \mathrm{CH}^{p}(S_1\times\cdots\times S_m)_{\mathbb Q}
 \longrightarrow
 H^{2p}(S_1\times\cdots\times S_m,\mathbb Q)
       \cap H^{p,p}(S_1\times\cdots\times S_m,\mathbb C)
\]
is surjective.  The intersection is taken after extension of scalars to
complex cohomology.
\end{assumption}

\begin{proposition}[Rational Hodge for products of geometric-K3 cubics]
\label{prop:cubic-products}
Assume Assumption~\ref{ass:k3-products}.  Let $X_1,\ldots,X_r$ be smooth
complex cubic fourfolds, with $r\geq 1$.  For each $i$, suppose there is a
smooth projective complex K3 surface $S_i$ and a Fourier--Mukai equivalence
\[
 \operatorname{Ku}(X_i)\simeq D^b(\operatorname{Coh}S_i),
\]
where $\operatorname{Ku}(X_i)$ is the Kuznetsov component and the K3 category
on the right is untwisted.  Then the rational Hodge conjecture holds in every
codimension on $X_1\times\cdots\times X_r$.  In particular, it holds on every
self-power of each such cubic fourfold.
\end{proposition}

\begin{proof}
We work with Chow motives over $\mathbb C$ with rational coefficients, using
the convention that $\mathbb Q(-1)$ has its Betti realization in degree two.
The algebraic and transcendental Chow--K\"unneth decomposition for a smooth
cubic fourfold, as recalled in \cite[\S5.2]{fuvial}, gives
\[
 h(X_i)\simeq
 \mathbb Q(0)\oplus\mathbb Q(-1)
 \oplus\mathbb Q(-2)^{\oplus a_i}
 \oplus h^4_{\mathrm{tr}}(X_i)
 \oplus\mathbb Q(-3)\oplus\mathbb Q(-4),
\]
where $a_i$ is the dimension of its rational algebraic middle cohomology.
Fu--Vial's theorem \cite[Theorem~3, equivalently Theorem~7.2]{fuvial}
applies to the stated Fourier--Mukai equivalence and gives an isomorphism
of rational Chow motives
\[
 h^4_{\mathrm{tr}}(X_i)(2)
       \simeq h^2_{\mathrm{tr}}(S_i)(1).
\]
Thus the non-Tate summand of $h(X_i)$ is isomorphic to
$h^2_{\mathrm{tr}}(S_i)(-1)$, itself an algebraic direct summand of a Tate
twist of $h(S_i)$.  These isomorphisms and their inverses are algebraic
correspondences, not merely isomorphisms of Hodge structures.

Tensor the displayed decompositions for all $i$.  Each summand of
$h(X_1\times\cdots\times X_r)$ is isomorphic to an algebraic direct
summand of a Tate twist of $h(\prod_{i\in I}S_i)$ for some subset
$I\subseteq\{1,\ldots,r\}$; for $I=\varnothing$ it is a Tate motive.
Assumption~\ref{ass:k3-products} applies to every nonempty such product.
Its conclusion passes to algebraic direct summands: if $\pi$ is the
corresponding algebraic projector and $\alpha$ is a Hodge class in its
image, first represent $\alpha$ by an algebraic cycle $z$ on the full
product and then use $\pi_*z$.  Tate twists only shift the degree and
codimension.  Finally transport these cycle classes through the inverse
algebraic correspondences and sum the components.  This represents every
rational Hodge class on the cubic product by an algebraic cycle.
\end{proof}

\begin{proposition}[Application to the released irrational cubics]
\label{prop:irrational-cubics}
Assume Assumption~\ref{ass:k3-products} and Theorem~1.1 of \cite{cubic}.
For every sufficiently large admissible Hassett discriminant $d$, a very
general smooth cubic fourfold $X$ in the Hassett divisor $\mathcal C_d$
is irrational and satisfies the rational Hodge conjecture on every
self-power.  More generally, the rational Hodge conjecture holds on every
finite product of cubics chosen in these very-general loci, even when the
discriminants differ.
\end{proposition}

\begin{proof}
The cited cubic theorem supplies irrationality and an exact complex-linear
equivalence of $\operatorname{Ku}(X)$ with the ordinary derived category
of a smooth projective K3 surface.  For this untwisted target the equivalence
is Fourier--Mukai, as recorded following Theorem~3 of \cite{fuvial}.
Proposition~\ref{prop:cubic-products} therefore applies to any finite list
of these cubics.
\end{proof}

\begin{remark}
Here admissible means $d>6$, $d\equiv 0$ or $2\pmod 6$, with $4\nmid d$,
$9\nmid d$, and no odd prime congruent to $2\pmod 3$ dividing $d$, as in
\cite{cubic}.  The discriminant threshold is ineffective, and ``very
general'' is interpreted separately in each $\mathcal C_d$.  The proposition
does not assert the Hodge conjecture for all cubic fourfolds, the full Hodge
conjecture, or stable irrationality.  The correspondence bridge is an
established theorem of Fu and Vial; the release supplies the K3-product and
irrationality inputs.  Novelty of the combined corollary is not established.
\end{remark}

\section{Idempotent lifting with necessarily infinite support}
\label{sec:adic-support}

Write $\mathbb Z_2$ for the ring of $2$-adic integers and
$\mathbb F_2$ for the field with two elements.  For a discrete group $G$,
$R[G]$ denotes its ordinary group ring: all its elements have finite group
support.  The completion below completes coefficients; it does not replace
$G$ by a profinite group.

\begin{assumption}[The two group-ring inputs]
\label{ass:group-rings}
The following statements from the release hold.
\begin{enumerate}
\item Theorem~1.1 of \cite{group2}: there exist a finitely presented
      torsion-free group $G$ with a finite two-dimensional classifying
      complex and $a,b,c\in\mathbb F_2[G]$ satisfying
      $ab=1$, $ac=0$, and $c\neq 0$.
\item Corollary~1.2(ii) of \cite{idempotent}: for every torsion-free
      discrete group $H$ and every commutative unital characteristic-zero
      domain $R$, the only idempotents of $R[H]$ are $0$ and $1$.
\end{enumerate}
The second statement concerns scalar ring elements, not arbitrary matrices.
\end{assumption}

\begin{proposition}[Completion creates an idempotent requiring infinite support]
\label{prop:adic-support}
Assume Assumption~\ref{ass:group-rings}, and take its group $G$ and elements
$a,b,c$.  Set
\[
 A=\mathbb Z_2[G],\qquad
 \widehat A=\varprojlim_{n\geq1}(\mathbb Z/2^n\mathbb Z)[G],
 \qquad e=1-ba\in\mathbb F_2[G].
\]
Then $e$ is a nontrivial idempotent and lifts to an idempotent of
$\widehat A$.  The ring $A$ has no nontrivial idempotents, and every
idempotent of $\widehat A$ reducing to $e$ has infinite support in $G$.
For every compatible sequence of idempotents
\[
 E_n\in(\mathbb Z/2^n\mathbb Z)[G],\qquad E_1=e,
\]
the finite numbers $|\operatorname{supp}E_n|$ tend to infinity as
$n\to\infty$.
\end{proposition}

\begin{proof}
First $ba\neq 1$, since otherwise
$c=(ba)c=b(ac)=0$.  As $ab=1$, we have $(ba)^2=b(ab)a=ba$, and hence
$e^2=e$.  The element $e$ is nonzero.  It is not $1$: if $ba=0$, then
$b=b(ab)=(ba)b=0$, contradicting $ab=1$.

Choose a finite-support integral lift $x_0\in\mathbb Z[G]\subset A$ of
$e$.  Define recursively
\[
 x_{j+1}=3x_j^2-2x_j^3,\qquad q_j=x_j^2-x_j.
\]
The following identities hold in any unital ring:
\[
 q_{j+1}=q_j^2(4q_j-3),\qquad
 x_{j+1}-x_j=(1-2x_j)q_j.
\]
Indeed, they are polynomial identities in the single element $x_j$ and
central integer scalars, so noncommutativity of the ambient ring causes
no issue.  Since $q_0\in2A$, induction gives
$q_j\in2^{2^j}A$.  The second identity makes $(x_j)$ a $2$-adic Cauchy
sequence.  Its limit $E\in\widehat A$ satisfies $E^2=E$ and reduces to
$e$ modulo $2$, so it is nontrivial.

For clarity, an element of $\widehat A$ can be identified with a
coefficient family $(u_g)_{g\in G}$ in $\mathbb Z_2$ such that, for every
$n$, only finitely many $u_g$ are nonzero modulo $2^n$.  Its support is
$\{g:u_g\neq0\}$.  Such an element has finite support exactly when it
lies in the natural embedded copy of $A$.

The second input, with $R=\mathbb Z_2$ and $H=G$, excludes every
nontrivial idempotent of $A$.  Thus any idempotent lift of $e$ in
$\widehat A$ has infinite support.  For a compatible sequence $(E_n)$,
its limit is such a lift.  Compatibility makes the sets
$\operatorname{supp}E_n$ nested, and their union is the infinite support
of the limit.  A nested sequence of finite sets with infinite union has
cardinalities tending to infinity, proving the last assertion.
\end{proof}

\begin{proposition}[Failure of direct finiteness after completion]
\label{prop:adic-direct-finiteness}
The first input of Assumption~\ref{ass:group-rings} alone implies that
$\widehat A$ is not directly finite: there exist $U,V\in\widehat A$
with $UV=1$ and $VU\neq1$.
\end{proposition}

\begin{proof}
Choose finite-support lifts $A_0,B_0\in A$ of $a,b$.  Then
$A_0B_0=1+2C$ for some $C\in A$.  The geometric series
$\sum_{j\geq0}(-2C)^j$ converges in $\widehat A$ and is a two-sided
inverse of $A_0B_0$.  Put
$U=A_0$ and $V=B_0(A_0B_0)^{-1}$.  Then $UV=1$, whereas
$VU$ reduces to $ba\neq1$ modulo $2$.
\end{proof}

\begin{remark}
The obstruction is finite support, not the existence of lifts at finite
precision.  No quantitative lower bound on support growth is asserted.
The lifting polynomials and the geometric-series argument are standard
completion methods; the conditional separation combines them with the
specified release inputs.  Its prior publication status is not established.
\end{remark}
